\section{Preliminaries}
\label{sec:preliminaries}

\subsection{Notation and Simplicial Mesh}
\label{subsec:notation}

Let $\Omega \subset \mathbb{R}^d$ ($d=2,3$) be a bounded Lipschitz polyhedral domain with boundary $\partial\Omega = \overline{\Gamma_D} \cup \overline{\Gamma_N}$, where $\Gamma_D \cap \Gamma_N = \emptyset$ and $|\Gamma_D| > 0$. Here $\Gamma_D$ denotes the Dirichlet boundary where displacements are prescribed, and $\Gamma_N$ represents the Neumann boundary subjected to surface tractions. Let $\mathbb{S}$ denote the space of real symmetric second-order tensors of order $d \times d$, and $\mathbb{I}$ be the second-order identity tensor.

For symmetric tensors $\boldsymbol{\sigma}, \boldsymbol{\tau} \in \mathbb{S}$, the Frobenius inner product and trace are denoted by $\boldsymbol{\sigma} : \boldsymbol{\tau} = \sum_{i,j=1}^d \sigma_{ij}\tau_{ij} = \operatorname{tr}(\boldsymbol{\sigma}\boldsymbol{\tau})$. For tensor field $\boldsymbol{\tau}$ and vector field $\boldsymbol{u}$, the divergence and symmetric gradient operators are defined by
\begin{equation}
(\operatorname{div}\boldsymbol{\tau})_i = \sum_{j=1}^d \frac{\partial \tau_{ij}}{\partial x_j}, \qquad 
\boldsymbol{\varepsilon}(\boldsymbol{u}) = \frac{1}{2}\left(\nabla\boldsymbol{u} + \nabla\boldsymbol{u}^{\mathsf T}\right).
\end{equation}

Let $\mathcal{T}_h$ be a shape-regular conforming simplicial mesh of $\Omega$ (triangles for $d=2$ and tetrahedra for $d=3$), characterized by mesh size $h = \max_{T\in\mathcal{T}_h} h_T$. We denote the set of vertices by $\mathcal{V}_h$, edges by $\mathcal{E}_h$, and $(d-1)$-dimensional faces by $\mathcal{F}_h$ (with $\mathcal{F}_h = \mathcal{E}_h$ in two dimensions). The face set is partitioned into internal faces $\mathcal{F}_h^i$, Dirichlet boundary faces $\mathcal{F}_h^D$, and Neumann boundary faces $\mathcal{F}_h^N$, such that $\mathcal{F}_h = \mathcal{F}_h^i \cup \mathcal{F}_h^D \cup \mathcal{F}_h^N$.

We adopt standard $L^2$ inner products and norms on $\Omega$, denoted by $(\cdot,\cdot)_\Omega$ and $\|\cdot\|_{0,\Omega}$. The symmetric tensor $H(\operatorname{div})$ space is defined as
\begin{equation}
H(\operatorname{div},\Omega;\mathbb{S}) = \left\{ \boldsymbol{\tau} \in L^2(\Omega;\mathbb{S}) : \operatorname{div}\boldsymbol{\tau} \in L^2(\Omega;\mathbb{R}^d) \right\},
\end{equation}
equipped with the graph norm $\|\boldsymbol{\tau}\|_{H(\operatorname{div})}^2 = \|\boldsymbol{\tau}\|_{0,\Omega}^2 + \|\operatorname{div}\boldsymbol{\tau}\|_{0,\Omega}^2$.

\subsection{Linear Elasticity and Hellinger--Reissner Variational Formulation}
\label{subsec:hellinger_reissner}

Consider the small-strain isotropic linear elasticity boundary value problem:
\begin{equation}
\label{eq:strong_elasticity}
\begin{aligned}
-\operatorname{div}\boldsymbol{\sigma} &= \boldsymbol{b} && \text{in } \Omega, \\
\mathcal{A}\boldsymbol{\sigma} &= \boldsymbol{\varepsilon}(\boldsymbol{u}) && \text{in } \Omega, \\
\boldsymbol{u} &= \boldsymbol{u}_D && \text{on } \Gamma_D, \\
\boldsymbol{\sigma}\boldsymbol{n} &= \boldsymbol{g} && \text{on } \Gamma_N,
\end{aligned}
\end{equation}
where $\boldsymbol{\sigma} \in \mathbb{S}$ is the symmetric Cauchy stress tensor, $\boldsymbol{u}$ is the displacement vector, $\boldsymbol{b}$ is the body force, $\boldsymbol{g}$ is the prescribed surface traction, and $\boldsymbol{n}$ is the unit outward normal vector on $\partial\Omega$. The fourth-order compliance tensor $\mathcal{A}$ acts on any symmetric tensor $\boldsymbol{\tau} \in \mathbb{S}$ as
\begin{equation}
\label{eq:compliance_tensor}
\mathcal{A}\boldsymbol{\tau} = \frac{1}{2\mu}\left( \boldsymbol{\tau} - \frac{\lambda}{2\mu + d\lambda} \operatorname{tr}(\boldsymbol{\tau})\mathbb{I} \right),
\end{equation}
with Lam{\'e} parameters $\lambda = \frac{E\nu}{(1+\nu)(1-2\nu)}$ and $\mu = \frac{E}{2(1+\nu)}$ for plane strain/3D, where $E$ is Young's modulus and $\nu$ is Poisson's ratio.

To handle the traction boundary condition $\boldsymbol{\sigma}\boldsymbol{n} = \boldsymbol{g}$ on $\Gamma_N$, we define the affine trial space and the homogeneous test space:
\begin{equation}
\boldsymbol{\Sigma}_g = \left\{ \boldsymbol{\tau} \in H(\operatorname{div},\Omega;\mathbb{S}) : \boldsymbol{\tau}\boldsymbol{n} = \boldsymbol{g} \text{ on } \Gamma_N \right\}, \qquad
\boldsymbol{\Sigma}_0 = \left\{ \boldsymbol{\tau} \in H(\operatorname{div},\Omega;\mathbb{S}) : \boldsymbol{\tau}\boldsymbol{n} = \boldsymbol{0} \text{ on } \Gamma_N \right\}.
\end{equation}
The displacement space is chosen as $\boldsymbol{V} = L^2(\Omega;\mathbb{R}^d)$. The two-field Hellinger--Reissner variational problem seeks $(\boldsymbol{\sigma},\boldsymbol{u}) \in \boldsymbol{\Sigma}_g \times \boldsymbol{V}$ such that
\begin{equation}
\label{eq:hr_variational}
\begin{aligned}
a(\boldsymbol{\sigma},\boldsymbol{\tau}) + b(\boldsymbol{\tau},\boldsymbol{u}) &= \langle \boldsymbol{u}_D, \boldsymbol{\tau}\boldsymbol{n} \rangle_{\Gamma_D}, && \forall\,\boldsymbol{\tau}\in\boldsymbol{\Sigma}_0, \\
b(\boldsymbol{\sigma},\boldsymbol{v}) &= -(\boldsymbol{b},\boldsymbol{v})_\Omega, && \forall\,\boldsymbol{v}\in\boldsymbol{V},
\end{aligned}
\end{equation}
where the bilinear forms and boundary duality pairing are given by
\begin{equation}
a(\boldsymbol{\sigma},\boldsymbol{\tau}) = \int_\Omega \mathcal{A}\boldsymbol{\sigma} : \boldsymbol{\tau} \,\mathrm{d}x, \qquad
b(\boldsymbol{\tau},\boldsymbol{v}) = \int_\Omega (\operatorname{div}\boldsymbol{\tau})\cdot\boldsymbol{v}\,\mathrm{d}x, \qquad
\langle \boldsymbol{u}_D, \boldsymbol{\tau}\boldsymbol{n} \rangle_{\Gamma_D} = \int_{\Gamma_D} \boldsymbol{u}_D \cdot (\boldsymbol{\tau}\boldsymbol{n})\,\mathrm{d}s.
\end{equation}
Because $\boldsymbol{u} \in L^2(\Omega;\mathbb{R}^d)$, no inter-element continuity is imposed on displacements. Meanwhile, as Poisson's ratio $\nu \to 0.5$, the compliance bilinear form $a(\cdot,\cdot)$ remains uniformly bounded, ensuring locking-free behavior in the continuous formulation~\cite{BoffiBrezziFortin2013}.

\subsection{Non-homogeneous Traction Lifting and Total Stress Formulation}
\label{subsec:traction_lifting}

In the mixed variational formulation~\eqref{eq:hr_variational}, the traction condition $\boldsymbol{\sigma}\boldsymbol{n} = \boldsymbol{g}$ acts as an essential boundary condition on the stress field. To enable discretization within a standard linear subspace, we introduce a design-independent traction lifting field $\boldsymbol{\sigma}_g \in \boldsymbol{\Sigma}_g$ satisfying $\boldsymbol{\sigma}_g\boldsymbol{n} = \boldsymbol{g}$ on $\Gamma_N$. The total stress is then decomposed into
\begin{equation}
\label{eq:stress_decomposition}
\boldsymbol{\sigma} = \boldsymbol{\sigma}_0 + \boldsymbol{\sigma}_g, \qquad \boldsymbol{\sigma}_0 \in \boldsymbol{\Sigma}_0.
\end{equation}
Substituting~\eqref{eq:stress_decomposition} into~\eqref{eq:hr_variational}, the problem is recast as: find $(\boldsymbol{\sigma}_0,\boldsymbol{u}) \in \boldsymbol{\Sigma}_0 \times \boldsymbol{V}$ such that
\begin{equation}
\label{eq:lifted_variational}
\begin{aligned}
a(\boldsymbol{\sigma}_0,\boldsymbol{\tau}) + b(\boldsymbol{\tau},\boldsymbol{u}) &= \langle \boldsymbol{u}_D, \boldsymbol{\tau}\boldsymbol{n} \rangle_{\Gamma_D} - a(\boldsymbol{\sigma}_g,\boldsymbol{\tau}), && \forall\,\boldsymbol{\tau}\in\boldsymbol{\Sigma}_0, \\
b(\boldsymbol{\sigma}_0,\boldsymbol{v}) &= -(\boldsymbol{b},\boldsymbol{v})_\Omega - b(\boldsymbol{\sigma}_g,\boldsymbol{v}), && \forall\,\boldsymbol{v}\in\boldsymbol{V}.
\end{aligned}
\end{equation}
Once $\boldsymbol{\sigma}_0$ is obtained, the total stress field $\boldsymbol{\sigma} = \boldsymbol{\sigma}_0 + \boldsymbol{\sigma}_g$ is reconstructed and subsequently utilized for objective evaluation and stress constraint enforcement throughout the topology optimization procedure.
